\documentclass[12pt]{report}

\usepackage{listings}
\usepackage{float}
\usepackage{booktabs}
\usepackage{amssymb}
\usepackage{subcaption}

\usepackage[dvipsnames, table, x11names, svgnames]{xcolor}
\usepackage{caption}
\usepackage{anyfontsize}
\usepackage{etoolbox}
\usepackage[italian]{babel}
\usepackage[pass]{geometry}
\usepackage{tcolorbox}
\usepackage{ragged2e}
\usepackage{listings}
\usepackage{ltablex}
\usepackage{environ}
\usepackage{titlesec}
\usepackage{tabularx}
\usetikzlibrary{positioning}
\usepackage[hidelinks]{hyperref}

\tcbuselibrary{theorems, listings, breakable, skins}

% Togliamo la fastidiosissima cosa che appena vado a capo mi mette lo skip
% come se fosse un paragrafo. Modifichiamo anche lo skip verticale.
\setlength{\parindent}{0pt}
\setlength{\parskip}{1ex}

% Modifichiamo lo stile del generico capitolo
\titleformat{\chapter}[display]
{\normalfont\huge\bfseries\linespread{1.3}\selectfont\raggedright}
{\chaptertitlename\ \thechapter}
{5pt}
{\Huge}

\titlespacing*{\chapter}{0pt}{0pt}{2pt}
\titlespacing*{\section}{0pt}{5pt}{3pt}
\titlespacing*{\subsection}{0pt}{5pt}{3pt}
\titlespacing*{\subsubsection}{0pt}{5pt}{3pt}

\newcommand{\generaFrontespizio}[4]{
    \begin{titlepage}
		
		\begin{center}
			
			\hbox{}
			\vfill
			
			{
				\color{BrickRed}
				\bfseries
				\fontsize{40pt}{40pt}\selectfont
				#1
				\par
			}
			
			\vspace{0.5cm}
			
			\vfill
			\hbox{}
			
			\begin{minipage}{1\textwidth}
				\fontsize{12pt}{14pt}\selectfont \textit{I seguenti appunti sono a cura di \textbf{Catalin Ceban} e si basano sul corso di "#1" erogato nell'A.A. #2 
				\ifstrequal{#4}{m}{ dal professore}{\ifstrequal{#4}{p}{ dai professori}{ dalla professoressa}} #3 per il C.d.L Informatica (33503)}
			\end{minipage}
			
			\vspace{0.5cm}
			{\large \today}
			
		\end{center}
		
	
	\end{titlepage}

	\newgeometry{
		includefoot,
		left=3.5cm,
		right=3.5cm,
		top=2.9cm,
		bottom=2cm
	}

	\tableofcontents
	\clearpage
	
}

\colorlet{coloreTeoremi}{DarkGoldenrod2}
\colorlet{coloreDefinizioni}{Coral3}
\colorlet{coloreProposizioni}{LightSalmon3}
\colorlet{coloreCorollari}{MediumOrchid3}
\colorlet{coloreOsservazioni}{Burlywood4}
\colorlet{coloreSfondoBox}{gray!5}

\newlength{\boxDistanzaVerticale}
\setlength{\boxDistanzaVerticale}{7pt}
\newlength{\boxDistanzaOrizzontale}
\setlength{\boxDistanzaOrizzontale}{10pt}
\newlength{\boxSpessoreRiga}
\setlength{\boxSpessoreRiga}{3pt}

\tcbset{
	stileBox/.style={
		blanker,
		borderline west={\boxSpessoreRiga}{0pt}{#1},
		left=\boxDistanzaOrizzontale,
		toptitle=\boxDistanzaVerticale, 
		bottom=\boxDistanzaVerticale,
		right=\boxDistanzaOrizzontale,
		pad after break=2mm,
		colback=coloreSfondoBox,
		colbacktitle=tcbcolback,
		coltitle=#1,
		fonttitle=\bfseries,
		enhanced,
		breakable
	}
}

\newtcbtheorem[number within=section]{teorema}{Teorema}{stileBox=coloreTeoremi}{th}

\newtcbtheorem[number within=section]{definizione}{Definizione}{stileBox=coloreDefinizioni}{def}

\newtcbtheorem[number within=section]{proposizione}{Proposizione}{stileBox=coloreProposizioni}{prop}

\newtcbtheorem[number within=section]{corollario}{Corollario}{stileBox=coloreCorollari}{cor}

\newtcbtheorem[number within=section]{osservazione}{Osservazione}{stileBox=coloreOsservazioni}{oss}

\newtcbtheorem[number within=section]{dimostrazione}{Dimostrazione}{stileBox=coloreOsservazioni}{oss}

\newtcbtheorem[number within=section]{algoritmo}{Spiegazione algoritmo}{stileBox=coloreProposizioni}{alg}

\newtcbtheorem[number within=section]{esercizio}{Esercizio}{stileBox=coloreTeoremi}{es}

\lstdefinestyle{stile-listings}{
	basicstyle=\ttfamily\small,
	inputencoding=utf8,
	extendedchars=true,
	literate=
	{à}{{\`a}}1
	{è}{{\`e}}1
	{é}{{\'e}}1
	{ì}{{\`i}}1
	{ò}{{\`ò}}1
	{ù}{{\`u}}1,
	keywordstyle=\color{blue!55!black}\bfseries,
	stringstyle=\color{Coral!55!black},
	numbers=left,
	numberstyle=\tiny\color{black!45},
	numbersep=8pt,
	backgroundcolor=\color{gray!15},
	frame=single,
	rulecolor=\color{gray!70!black},
	framesep=6pt,
	framextopmargin=4pt,
	framexbottommargin=4pt,
	xleftmargin=7pt,
	xrightmargin=7pt,
	breaklines=true,
	breakatwhitespace=false,
	showstringspaces=false,
	tabsize=2,
	captionpos=b,
	aboveskip=15pt,
	belowskip=8pt
}

\newcolumntype{C}{>{\Centering\arraybackslash}X}

\NewEnviron{tabella}[2][]{%
	\par\vspace{0.5em}%
	\begingroup
	\centering
	\renewcommand{\arraystretch}{1.5}%
	\begin{tabularx}{\textwidth}{#2}
		\hline
		\rowcolor{BrickRed!60!white}
		\BODY
		\hline
	\end{tabularx}%
	\if\relax\detokenize{#1}\relax
	% Nessuna caption fornita
	\else
	\vspace{0.5em}%
	\captionof{table}{#1}%
	\fi
	\par
	\endgroup
	\vspace{0.5em}%
}



\lstset{style=stile-listings}

\begin{document}
	
	\generaFrontespizio{Algoritmi 2}{2026}{Paul Joseph Wollan}{m}
	
	
	% --- NEW PAGE (1) ---
	
	\chapter{Introduzione e Teoria dei Grafi}
	
	In questa sezione introduttiva delineiamo gli argomenti principali che saranno trattati nel corso, ponendo le basi per lo studio strutturato degli algoritmi avanzati.
	
	\section{Argomenti principali}
	Il corso si articolerà principalmente attorno ai seguenti nuclei tematici:
	\begin{enumerate}
		\item \textbf{Grafi e algoritmi sui grafi}: studio delle strutture dati fondamentali per la rappresentazione di relazioni e delle relative tecniche di visita ed elaborazione.
		\item \textbf{Algoritmi greedy}: strategie di scelta ottima locale per la risoluzione di problemi globali (ad esempio, la ricerca di un albero di copertura di peso minimo connettendo vari nodi).
		\item \textbf{Divide et impera}: paradigma algoritmico basato sull'approccio \textit{top-down}, che consiste nel suddividere un problema in sottoproblemi più piccoli, risolverli ricorsivamente e combinarne le soluzioni.
		\item \textbf{Programmazione dinamica}: approccio di tipo \textit{bottom-up} per la risoluzione di problemi caratterizzati da sottoproblemi sovrapposti, memorizzando i risultati parziali per evitare ricalcoli superflui.
	\end{enumerate}
	
	\begin{osservazione}{Albero}{}
		Ricordiamo che un \textbf{albero} è definito formalmente come un grafo connesso e aciclico.
	\end{osservazione}
	
	\section{Definizioni preliminari sui grafi}
	
	Un \textbf{grafo} è una coppia $G = (V, E)$, dove:
	\begin{itemize}
		\item $V$ rappresenta l'insieme dei vertici o nodi.
		\item $E$ rappresenta l'insieme degli archi o lati.
	\end{itemize}
	Per comodità di notazione, nel corso del manuale denoteremo con $n = |V|$ il numero dei vertici e con $m = |E|$ il numero degli archi.
	
	Ogni arco è una coppia di punti che indica i due vertici collegati. A titolo esemplificativo, consideriamo la seguente rappresentazione grafica:
	
	\begin{center}
		\begin{tikzpicture}[scale=1.1, every node/.style={circle, draw, inner sep=1.5pt}]
			\node (A) at (0, 1) {};
			\node (B) at (2, 1) {};
			\node (C) at (1, 0) {};
			\node (D) at (1, 1.8) {};
			
			\draw (A) -- (B);
			\draw (A) -- (C);
			\draw (B) -- (C);
			\draw (A) -- (D);
			\draw (B) -- (D);
		\end{tikzpicture}
	\end{center}
	
	Riguardo agli archi, potremmo avere anche un orientamento: se ogni arco possiede una direzione definita, diciamo che il grafo è \textbf{diretto}. 
	
	\begin{center}
		\begin{tikzpicture}[scale=1.1, every node/.style={circle, draw, inner sep=1.5pt}]
			\node (A) at (0, 0) {};
			\node (B) at (1.5, 1.2) {};
			\node (C) at (3, 0) {};
			
			\draw[->, >=stealth] (A) -- (B);
			\draw[->, >=stealth] (B) -- (C);
			\draw[->, >=stealth] (C) -- (A);
			\draw[<->, >=stealth, bend left=30] (A) -- (C);
		\end{tikzpicture}
	\end{center}
	
	È importante notare che possiamo avere anche due archi sugli stessi nodi ma con orientamento diverso (come mostrato nel collegamento bidirezionale o nei doppi archi orientati tra i medesimi vertici).
	
	Sia $f$ un arco che collega due vertici $x$ e $y$, formalmente rappresentabile come $x \longleftarrow f \longrightarrow y$. Diciamo che:
	\begin{itemize}
		\item L'arco $f$ è \textbf{incidente} ai vertici $x$ e $y$.
		\item Il vertice $x$ è \textbf{adiacente} al vertice $y$.
	\end{itemize}
	
	\begin{definizione}{Grado di un vertice}{}
		Il \textbf{grado} di un vertice $x$, denotato con $\deg(x)$, è il numero di archi incidenti in $x$.
	\end{definizione}
	
	Ad esempio, considerando un nodo centrale collegato a quattro vertici periferici:
	
	\begin{center}
		\begin{tikzpicture}[scale=1.1, every node/.style={circle, draw, inner sep=1.5pt}]
			\node (C) at (1, 1) {};
			\node (1) at (0, 2) {};
			\node (2) at (2, 2) {};
			\node (3) at (0.5, 0) {};
			\node (4) at (1.5, 0) {};
			
			\draw (C) -- (1);
			\draw (C) -- (2);
			\draw (C) -- (3);
			\draw (C) -- (4);
			
			\draw (1) -- (2);
			\draw (3) -- (4);
		\end{tikzpicture}
	\end{center}
	
	Per il nodo centrale $x$ contrassegnato, avremo chiaramente $\deg(x) = 4$, poiché vi incidono esattamente quattro archi distinti.
	% --- NEW PAGE (2) ---
	
	\chapter{Introduzione ai grafi e passeggiate euleriane}
	
	\section{Il problema dei ponti di K\"onigsberg}
	
	Il problema classico da cui nasce la teoria dei grafi è il famoso problema dei ponti di K\"onigsberg:
	\begin{center}
		\begin{tikzpicture}[scale=0.8]
			% Fiumi e terre schematici
			\draw[thick] (0,1) to[out=10,in=170] (6,1);
			\draw[thick] (0,-1) to[out=-10,in=-170] (6,-1);
			\draw[thick] (2.5,-1.3) to[out=80,in=280] (2.5,1.3);
			\draw[thick] (3.5,-1.3) to[out=100,in=260] (3.5,1.3);
			
			% Terre (isole/rive)
			\draw[thick, fill=gray!20] (1,0) ellipse (0.8cm and 0.5cm);
			\draw[thick, fill=gray!20] (5,0) ellipse (0.8cm and 0.5cm);
			\draw[thick, fill=gray!20] (3,1.5) rectangle (4,1.8);
			\draw[thick, fill=gray!20] (3,-1.8) rectangle (4,-1.5);
			
			% Ponti
			\draw[double,thick] (1.3,0.5) -- (1.3,1.5);
			\draw[double,thick] (1.7,0.5) -- (1.7,1.5);
			
			\draw[double,thick] (1.3,-0.5) -- (1.3,-1.5);
			\draw[double,thick] (1.7,-0.5) -- (1.7,-1.5);
			
			\draw[double,thick] (1.8,0) -- (4.2,0);
			
			\draw[double,thick] (4.7,0.5) -- (4.7,1.5);
			\draw[double,thick] (5.1,0.5) -- (5.1,1.5);
			
			\draw[double,thick] (4.7,-0.5) -- (4.7,-1.5);
			\draw[double,thick] (5.1,-0.5) -- (5.1,-1.5);
		\end{tikzpicture}
	\end{center}
	
	Ci si chiede: \textit{è possibile attraversare tutti i ponti una sola volta e arrivare al punto di partenza?}
	
	Trduciamo il problema in un problema di grafi:
	\begin{center}
		\begin{tikzpicture}[scale=0.7]
			% Semplice schema di traduzione
			\node at (0,0) {
				\begin{tikzpicture}[scale=0.5]
					\draw[thick, fill=gray!10] (0,0) ellipse (0.6cm and 0.4cm);
					\draw[thick, fill=gray!10] (3,0) ellipse (0.6cm and 0.4cm);
					\draw[double,thick] (0.6,0.2) -- (2.4,0.2);
					\draw[double,thick] (0.6,-0.2) -- (2.4,-0.2);
				\end{tikzpicture}
			};
			\node at (2.2,0) {$\Rightarrow$};
			
			% Grafo corrispondente
			\begin{scope}[shift={(4,-1.5)}, vertex/.style={circle, draw, inner sep=1.5pt}]
				\node[vertex, label=above:{\small 1}] (v1) at (1,3) {};
				\node[vertex, label=above:{\small 2}] (v2) at (2,2.2) {};
				\node[vertex, label=above:{\small 3}] (v3) at (3,2) {};
				\node[vertex, label=right:{\small 4}] (v4) at (5,1.8) {};
				\node[vertex, label=below:{\small 5}] (v5) at (3,1) {};
				\node[vertex, label=right:{\small 6}] (v6) at (5,0.8) {};
				\node[vertex, label=right:{\small 7}] (v7) at (5.5,0.2) {};
				\node[vertex, label=left:{\small 8}] (v8) at (0.5,0.5) {};
				\node[vertex, label=below:{\small 9}] (v9) at (3,0) {};
				\node[vertex, label=below:{\small 10}] (v10) at (5,-0.8) {};
				\node[vertex, label=below:{\small 11}] (v11) at (2,-1) {};
				
				% Archi
				\draw (v1) -- node[above, green!50!black, font=\small] {$a$} (v2);
				\draw (v2) -- node[above, green!50!black, font=\small] {$b$} (v3);
				\draw (v3) -- node[above, green!50!black, font=\small] {$c$} (v4);
				\draw (v1) -- node[left, green!50!black, font=\small] {$d$} (v8);
				\draw (v2) -- node[above, green!50!black, font=\small] {$e$} (v5);
				\draw (v4) -- node[right, green!50!black, font=\small] {$f$} (v7);
				\draw (v5) -- node[above, green!50!black, font=\small] {$n$} (v9);
				\draw (v5) -- node[below, green!50!black, font=\small] {$g$} (v9); % doppio arco approssimato
				\draw (v6) -- node[above, green!50!black, font=\small] {$i$} (v7);
				\draw (v8) -- node[below, green!50!black, font=\small] {$j$} (v9);
				\draw (v9) -- node[below, green!50!black, font=\small] {$m$} (v11);
			\end{scope}
		\end{tikzpicture}
	\end{center}
	
	\begin{definizione}{Passeggiata in un grafo}{}
		Una \textbf{passeggiata} in un grafo è una sequenza $v_0, e_1, v_1, e_2, v_2, \dots, v_{k-1}, e_k, v_k$ che alterna vertici e archi.
	\end{definizione}
	
	\begin{definizione}{Passeggiata chiusa}{}
		Una passeggiata si dice \textbf{chiusa} se $v_0 = v_k$, cioè il punto di inizio coincide con il punto di arrivo.
	\end{definizione}
	
	Qualsiasi grafo con almeno un nodo $x$ con $\deg(x)$ dispari \textbf{non} ammette una passeggiata chiusa che utilizza tutti gli archi una sola volta.
	In questo caso, dunque: \textbf{no}, non è risolvibile il problema.
	
	E se fossero tutti pari? \textbf{Non è detto}.
	
	\begin{center}
		\begin{tikzpicture}[scale=0.6]
			% Due triangoli disconnessi
			\node[vertex/.style={circle, draw, inner sep=1.5pt}] (a1) at (0,0) {};
			\node[vertex/.style={circle, draw, inner sep=1.5pt}] (a2) at (1,0) {};
			\node[vertex/.style={circle, draw, inner sep=1.5pt}] (a3) at (0.5,0.8) {};
			\draw (a1) -- (a2) -- (a3) -- (a1);
			
			\node[vertex/.style={circle, draw, inner sep=1.5pt}] (b1) at (0,-1.5) {};
			\node[vertex/.style={circle, draw, inner sep=1.5pt}] (b2) at (1,-1.5) {};
			\node[vertex/.style={circle, draw, inner sep=1.5pt}] (b3) at (0.5,-2.3) {};
			\draw (b1) -- (b2) -- (b3) -- (b1);
		\end{tikzpicture}
	\end{center}
	Tutti $\text{deg} = 2$, ma non essendo connesso non è possibile.
	
	\begin{definizione}{Grafo connesso}{}
		Un grafo si dice \textbf{connesso} se per ogni coppia di vertici esiste una passeggiata da $x$ a $y$.
	\end{definizione}
	% --- NEW PAGE (3) ---
	
	\chapter{Rappresentazione dei grafi}
	
	\begin{teorema}{Caratterizzazione di Eulero}{}
		Esiste una passeggiata chiusa che contiene ogni arco una sola volta se e solo se:
		\begin{enumerate}
			\item Ogni vertice ha grado pari.
			\item Il grafo è connesso.
		\end{enumerate}
	\end{teorema}
	
	\section{Rappresentazione fisica dei grafi}
	
	La rappresentazione informatica dei grafi in memoria è fondamentale per determinare l'efficienza degli algoritmi che operano su di essi. Esistono principalmente due modalità standard per rappresentare un grafo: le matrici di adiacenza e le liste di adiacenza.
	
	\subsection{Matrici di adiacenza}
	Si stabilisce una matrice $n \times n$ dove $n = \#V$, e dove:
	\[
	m_{i,j} = \begin{cases}
		1 & \text{se } i \text{ è adiacente a } j \\
		0 & \text{altrimenti}
	\end{cases}
	\]
	
	\textbf{Esempio:}
	\begin{center}
		\begin{minipage}{0.45\textwidth}
			\centering
			\begin{tikzpicture}[scale=0.9, every node/.style={circle,draw,inner sep=2pt}]
				\node (1) at (0,1) {$1$};
				\node (2) at (1.5,2) {$2$};
				\node (3) at (1.5,0) {$3$};
				\node (4) at (3,1) {$4$};
				
				\draw (1) -- (2);
				\draw (1) -- (3);
				\draw (2) -- (3);
				\draw (2) -- (4);
				\draw (3) -- (4);
			\end{tikzpicture}
		\end{minipage}
		\begin{minipage}{0.45\textwidth}
			\centering
			\[
			\begin{matrix}
				& \begin{matrix} 1 & 2 & 3 & 4 \end{matrix} \\
				\begin{matrix} 1 \\ 2 \\ 3 \\ 4 \end{matrix} &
				\begin{bmatrix}
					0 & 1 & 1 & 0 \\
					1 & 0 & 1 & 1 \\
					1 & 1 & 0 & 1 \\
					0 & 1 & 1 & 0
				\end{bmatrix}
			\end{matrix}
			\]
		\end{minipage}
	\end{center}
	
	Se il grafo non è diretto, la matrice sarà simmetrica rispetto alla diagonale.
	
	\textbf{Costo computazionale:}
	\begin{itemize}
		\item Visita del grafo e memorizzazione: $\Theta(n^2)$.
		\item Per verificare se $i$ è adiacente a $j$: $\Theta(1)$.
		\item Verificare $\deg(i)$: $\Theta(n)$ (occorre scorrere la riga $i$-esima e contare quanti $1$ ci sono).
	\end{itemize}
	
	\subsection{Liste di adiacenza}
	
	Una lista di adiacenza è composta da un vettore di teste di liste concatenate. Ciascuna lista contiene i nodi adiacenti al rispettivo vertice.
	
	\begin{center}
		\begin{tabular}{c c}
			$\begin{bmatrix}
				1 \\ 2 \\ 3 \\ 4
			\end{bmatrix}$
			&
			$\begin{aligned}
				&\rightarrow [2, 3] \\
				&\rightarrow [1, 3, 4] \\
				&\rightarrow [1, 2, 4] \\
				&\rightarrow [2, 3]
			\end{aligned}$
		\end{tabular}
	\end{center}
	
	Il motivo per cui lo spazio complessivo dipende dal numero di archi è che ogni arco viene contato due volte, una per ogni nodo coinvolto nel grafo non orientato.
	
	\textbf{Costo computazionale (C.C.):}
	\begin{itemize}
		\item Per verificare se $i$ è adiacente a $j$: $\mathcal{O}(\deg(i))$.
		\item Per visitare e memorizzare: $\mathcal{O}\left(n + \sum_{v} \deg(v)\right) = \mathcal{O}(n + 2m) = \mathcal{O}(n + m)$, dove $n = \#V$ e $m = \#E$.
	\end{itemize}
	
	L'occupazione di memoria $\mathcal{O}(n)$ è dovuta chiaramente al fatto che deve esserci un array con un elemento per ogni nodo. Ogni nodo poi dovrà avere un puntatore ad una lista concatenata che rappresenta i suoi adiacenti.
	% --- NEW PAGE (4) ---
	
	\chapter{Passeggiate e Cammini nei Grafi}
	
	Ci poniamo ora la seguente domanda: \textit{come si estendono questi concetti nel caso dei grafi diretti?}
	
	\begin{center}
		\begin{tikzpicture}[->, >=stealth, thick, node distance=2cm]
			\node[circle, draw, inner sep=2pt] (1) {$1$};
			\node[circle, draw, inner sep=2pt] (2) [below right of=1] {$2$};
			\node[circle, draw, inner sep=2pt] (3) [above left of=2] {$3$};
			
			\draw (1) -- (2);
			\draw (2) -- (3);
			\draw (3) -- (1);
		\end{tikzpicture}
		\hspace{2cm}
		$\begin{bmatrix} 1 \\ 2 \\ 3 \end{bmatrix} \longrightarrow \begin{bmatrix} 2 \\ 3 \\ 1 \end{bmatrix}$
	\end{center}
	
	C'è un altro modo per memorizzare questi concetti. Consideriamo un grafo non orientato con i vertici $1, 2, 3, 4, 5$ collegati da vari archi etichettati.
	
	\begin{center}
		\begin{tikzpicture}[thick, node distance=1.5cm]
			\node[circle, draw, inner sep=2pt] (1) at (0,0) {$1$};
			\node[circle, draw, inner sep=2pt] (2) at (2,0) {$2$};
			\node[circle, draw, inner sep=2pt] (3) at (3,1.5) {$3$};
			\node[circle, draw, inner sep=2pt] (4) at (1,1.5) {$4$};
			\node[circle, draw, inner sep=2pt] (5) at (4,0) {$5$};
			
			\draw (1) -- node[below] {$a$} (2);
			\draw (2) -- node[below] {$b$} (5);
			\draw (2) -- node[right] {$c$} (3);
			\draw (3) -- node[above] {$d$} (4);
			\draw (4) -- node[left] {$e$} (1);
			\draw (2) -- (4);
		\end{tikzpicture}
	\end{center}
	
	Una passeggiata all'interno del grafo potrebbe essere la seguente sequenza:
	\[
	1 \xrightarrow{a} 2 \xrightarrow{c} 3 \xrightarrow{d} 4 \xrightarrow{e} 2 \xrightarrow{b} 5
	\]
	
	\begin{definizione}{Cammino}{}
		Un \textbf{cammino} in un grafo è una passeggiata che non ripete i vertici.
	\end{definizione}
	
	La passeggiata vista prima \textbf{non} è un cammino da $1$ a $5$, ma un possibile cammino tra $1$ e $5$ potrebbe essere:
	\[
	1 \xrightarrow{a} 2 \xrightarrow{b} 5
	\]
	
	\begin{proposizione}{Esistenza di cammini e passeggiate}{}
		Esiste una passeggiata da $x$ a $y$ se e solo se esiste un cammino da $x$ a $y$.
	\end{proposizione}
	
	\begin{dimostrazione}{}{}
		Chiaramente, se esiste un cammino allora sicuramente esiste anche una passeggiata, dato che ogni cammino è un particolare tipo di passeggiata.\\
		Viceversa, supponiamo di avere una passeggiata $W = (v_0, v_1, \dots, v_K)$. Abbiamo due possibilità:
		\begin{enumerate}
			\item Se tutti i vertici sono distinti, allora la passeggiata è già un cammino.
			\item Se invece abbiamo almeno un indice ripetuto, cioè esistono $i, j$ tali che $v_i = v_j$ con $i < j$.
		\end{enumerate}
		Questo significa che all'interno della passeggiata è presente un sottociclo:
		\[
		W_1 = (v_i, v_{i+1}, \dots, v_{j-1}, v_j)
		\]
		ma sapendo che $v_i = v_j$, la struttura si chiude. Possiamo eliminare $W_1$ dalla passeggiata iniziale mantenendo però il nodo in cui parte la diramazione (in questo caso $v_i$). La passeggiata iniziale può essere dunque riscritta come:
		\[
		W = (v_1, \dots, v_i, \underbrace{v_{i+1}, \dots, v_{j-1}, v_j}_{W_1}, v_{j+1}, \dots, v_K)
		\]
		Eliminando $W_1$ otteniamo un cammino più corto, e reiterando il procedimento rimuoviamo tutti i cicli ottenendo infine un cammino semplice.
	\end{dimostrazione}
	
	\begin{definizione}{Ciclo}{}
		Un \textbf{ciclo} è un grafo connesso in cui ogni vertice ha grado $2$.
	\end{definizione}
	
	\begin{center}
		\begin{tikzpicture}[thick, scale=0.8]
			\node[circle, draw, inner sep=2pt] (1) at (0,2) {};
			\node[circle, draw, inner sep=2pt] (2) at (1.5,1.5) {};
			\node[circle, draw, inner sep=2pt] (3) at (1,0) {};
			\node[circle, draw, inner sep=2pt] (4) at (-1,0) {};
			\node[circle, draw, inner sep=2pt] (5) at (-1.5,1.5) {};
			
			\draw (1) -- (2);
			\draw (2) -- (3);
			\draw (3) -- (4);
			\draw (4) -- (5);
			\draw (5) -- (1);
		\end{tikzpicture}
	\end{center}
	% --- NEW PAGE (5) ---
	
	\chapter{Ricerca di cicli nei grafi}
	
	Infine la passeggiata diventa:
	\[
	W = (v_1, v_2, \dots, v_i, v_{j+1}, \dots, v_k)
	\]
	che, non contenendo nodi ripetuti, è un cammino.
	
	\begin{definizione}{Problema del ciclo}{}
		Dato un grafo $G$ dove per ogni vertice $x$ si ha $\deg(x) \geq 2$, trovare un ciclo nel grafo.
	\end{definizione}
	
	Prociamo prima graficamente:
	
	\begin{center}
		\begin{tikzpicture}[scale=0.9, every node/.style={circle, draw, thick, inner sep=1.5pt}]
			% Vertici del cubo/grafo
			\node (v1) at (0,0) {};
			\node (v2) at (2,0) {};
			\node (v3) at (2.5,1) {};
			\node (v4) at (0.5,1) {};
			\node (v5) at (0,1.8) {};
			\node (v6) at (2,1.8) {};
			\node (v7) at (2.5,2.8) {};
			\node (v8) at (0.5,2.8) {};
			\node (center) at (1.25, 1.4) {};
			
			% Archi
			\draw (v1) -- (v2);
			\draw (v1) -- (v4);
			\draw (v2) -- (v3);
			\draw (v4) -- (v3);
			\draw (v5) -- (v6);
			\draw (v5) -- (v8);
			\draw (v6) -- (v7);
			\draw (v8) -- (v7);
			\draw (v1) -- (v5);
			\draw (v2) -- (v6);
			\draw (v3) -- (v7);
			\draw (v4) -- (v8);
			\draw (center) -- (v3);
			\draw (center) -- (v5);
			\draw (center) -- (v2);
		\end{tikzpicture}
	\end{center}
	
	L'idea iniziale consiste nel:
	\begin{itemize}
		\item Prendere un punto di partenza.
		\item Prendere un arco e andare al nodo adiacente.
		\item Memorizzare i nodi visitati finora.
		\item Controllare se il prossimo nodo è stato già visitato.
	\end{itemize}
	
	Presentiamo di seguito l'algoritmo per la ricerca del ciclo all'interno del grafo $G$.
	
	\begin{algoritmo}{Algoritmo di ricerca del ciclo}{ciclo}
		Il procedimento riceve in input un grafo $G$ in cui ogni vertice ha grado almeno $2$. L'obiettivo è individuare un ciclo sfruttando una lista di supporto per tracciare i nodi visitati.
		\begin{lstlisting}[style=stile-listings]
			Ciclo(G):
			x = un vertice qualsiasi
			creo una LISTA PUNTATA vuota (L)
			append(x)
			teniamo traccia di current(x) e next = x.neighbour[0]
			while next not in L:
			L.append(next)
			if next.neighb[0] == current:
			u = next.neighb[1]
			else:
			u = next.neighb[0]
			current = next
			next = u
			while L[0] != next:
			L.remove(0)
			return L
		\end{lstlisting}
		\smallbreak
		\textbf{Note sull'implementazione e complessità:}
		\begin{itemize}
			\item La ricerca del nodo all'interno della lista $L$ (`while next not in L`) richiede normalmente $O(n)$, ma tale operazione può essere gestita facilmente definendo un vettore di supporto `vis = [0] * n` che si aggiorna ponendo `vis[next] = 1` ad ogni iterazione, riducendo il costo di controllo a $O(1)$.
			\item Se invece si decidesse di utilizzare una matrice di adiacenza, il controllo costerebbe comunque $O(n)$, portando il costo computazionale complessivo a $O(n^2)$.
		\end{itemize}
	\end{algoritmo}
	% --- NEW PAGE (6) ---
	
	\chapter{Ricerca di cicli nei grafi}
	
	L'idea di base dell'algoritmo risolutivo per la ricerca di un ciclo all'interno di un grafo (supportata dall'ipotesi forte che ogni vertice abbia grado $\deg(x) \ge 2$) si articola nei seguenti punti:
	\begin{enumerate}
		\item Si sceglie un nodo di partenza $x$ e lo si memorizza in una struttura dati che permetta di tenere traccia dei nodi visitati.
		\item Si prende un qualsiasi nodo adiacente, controllando categoricamente che non sia già stato visitato (per la prima iterazione è impossibile, ma in generale...).
	\end{enumerate}
	
	Siamo certi che esista un ciclo proprio grazie all'ipotesi forte di $\deg(x) \ge 2$ per ogni $x$. 
	Proviamo a formalizzare il tutto sotto forma di codice.
	
	\begin{lstlisting}[style=stile-listings, caption={Algoritmo per la ricerca di un ciclo in un grafo}]
		def trovaCiclo(G):
		L = lista vuota
		x = vertice qualsiasi
		L.append(x)
		Current = x
		next = x.neighbr[0]
		
		while (next not in L):
		L.append(next)
		if next.neighbr[0] == current:
		u = next.neighbr[1]
		else:
		u = next.neighbr[0]
		current = next
		next = u
		
		while L[0] != next:
		L.remove(0)
		
		return L
	\end{lstlisting}
	
	\begin{algoritmo}{Analisi e spiegazione dell'algoritmo}{}
		Il codice presentato implementa la ricerca di un ciclo all'interno del grafo $G$. Analizziamo dettagliatamente le singole componenti e le istruzioni evidenziate:
		
		\begin{itemize}
			\item \textbf{Inizializzazione:} Viene creata una lista vuota $L$ per tracciare il percorso e viene scelto un vertice iniziale $x$ aggiunto alla lista. Si impostano i puntatori iniziali \texttt{current} e \texttt{next}. Nota importante: l'operazione di controllo di presenza all'interno di una lista non ordinata (\texttt{next not in L}) richiede tempo $O(n)$, rendendo questa versione iniziale inefficiente.
			\item \textbf{Esplorazione (Righe 7-14):} Finché il nodo successivo \texttt{next} non è già presente in $L$\textsuperscript{1}, il nodo viene aggiunto alla lista\textsuperscript{2}. Si controlla quale dei vicini corrisponde al nodo corrente per evitare di tornare indietro sui propri passi: se il primo vicino è il nodo corrente, si sceglie il secondo, altrimenti si prende tranquillamente il primo dei vicini.
			\item \textbf{Aggiornamento dei puntatori (Riga 12-13):} Vengono preparati i puntatori per la successiva iterazione\textsuperscript{3} aggiornando \texttt{current} a \texttt{next} e \texttt{next} a $u$. Questo ciclo viene eseguito al più $n$ volte, portando a una complessità temporale complessiva di $T(n) = O(n^2)$.
			\item \textbf{Pulizia della struttura (Righe 16-17):} Una volta individuato il ciclo (quando \texttt{next} è già presente in $L$, il che significa che siamo in qualche modo tornati sui nostri passi), si rimuovono dalla lista $L$ tutti gli elementi iniziali che non fanno parte del ciclo, finché il primo elemento di $L$ non coincide con il nodo \texttt{next}.
		\end{itemize}
	\end{algoritmo}
	% --- NEW PAGE (7) ---
	
	\chapter{Rilevamento di cicli nei grafi}
	
	Quando raggiungiamo un ciclo, la variabile \texttt{next} contiene l'identificativo del vertice su cui il ciclo si è chiuso. Tuttavia, la lista $L$ contiene tutto il cammino percorso fino a quel momento, quindi dobbiamo eliminare tutto ciò che non fa parte del ciclo.
	
	Controlliamo se $L[0] \neq \text{next}$: sapendo che \texttt{next} è l'ultimo nodo del ciclo, significa che esso è anche il primo nodo del ciclo stesso. Dunque, il primo valore della lista dovrà essere proprio \texttt{next}, in quanto è il punto da cui tutto parte.
	
	\begin{figure}[htbp]
		\centering
		\begin{tikzpicture}[scale=1.1, every node/.style={circle, draw, minimum size=6mm}]
			\node (1) at (0, 1) {$1$};
			\node (2) at (1, 1) {$2$};
			\node (3) at (2, 1) {$3$};
			\node (4) at (2, 0) {$4$};
			\node (5) at (0, 0) {$5$};
			
			\draw (1) -- (2);
			\draw (2) -- (3);
			\draw (3) -- (4);
			\draw (4) -- (5);
			\draw (5) -- (3);
		\end{tikzpicture}
		\caption{Esempio grafico di un grafo con un ciclo.}
	\end{figure}
	
	Dopo il primo ciclo, abbiamo $L = [1, 2, 3, 4, 5]$, $\text{current} = 5$ e $\text{next} = 3$.
	Inizia il secondo ciclo di pulizia della lista:
	\begin{itemize}
		\item $L[0] \neq 3$? Sì, rimuoviamo $1$ e otteniamo $L = [2, 3, 4, 5]$.
		\item $L[0] \neq 3$? Sì, rimuoviamo $2$ e otteniamo $L = [3, 4, 5]$.
		\item $L[0] \neq 3$? No, il primo elemento è $3$: ciclo trovato, fine.
	\end{itemize}
	
	Tuttavia, come già detto, il contro di questa versione è che il controllo di presenza di un nodo nella lista può essere dispendioso. Possiamo infatti mitigare quel costo con l'uso di una lista \texttt{vis} che mappa i nodi visitati.
	
	\begin{algoritmo}{Algoritmo di ricerca del ciclo con vettore dei visitati}{}
		Viene presentata di seguito l'implementazione ottimizzata per la ricerca di un ciclo all'interno di un grafo $G$, utilizzando un vettore di supporto per tracciare i nodi visitati in tempo costante $\mathcal{O}(1)$.
		
		\begin{lstlisting}[style=stile-listings]
			def trovaCiclo(G):
			n = len(G.V)
			L = lista()
			vis = [0] * n
			
			x = vertice_qualsiasi()
			L.append(x)  # Aggiungiamo alla cronologia
			vis[x] = 1   # Segniamo come visitato
			
			current = x
			next = x.neighbr[0]
			
			while vis[next] == 0:  # Controllo in O(1)
			L.append(next)
			vis[next] = 1
			
			if next.neighbr[0] == current:
			u = next.neighbr[1]
			else:
			u = next.neighbr[0]
			
			current = next
			next = u
			
			# Rimozione dei nodi esterni al ciclo (stessa logica vista in precedenza)
			while L[0] != next:
			L.pop(0)
			
			return L
		\end{lstlisting}
	\end{algoritmo}
	% --- NEW PAGE (8) ---
	
	\chapter{Analisi delle versioni}
	
	Analizzando le due versioni, vediamo come:
	
	\begin{itemize}
		\item \textbf{Versione 1 (V1)}:
		\begin{itemize}
			\item Costo temporale: $O(n^2)$
			\item Costo spaziale: $O(n)$, ottenuto mediante l'utilizzo di una lista $L$ per tenere traccia della cronologia.
		\end{itemize}
		
		\item \textbf{Versione 2 (V2)}:
		\begin{itemize}
			\item Costo temporale: $O(n)$
			\item Costo spaziale: $O(n) + O(n) = O(2n) = O(n)$, dove i contributi spaziali sono dati rispettivamente dalla struttura $L$ e dall'insieme di supporto per le visite $Vis$. Poiché il risultato si attesta sullo stesso ordine di grandezza, la complessità spaziale risulta essere ottima.
		\end{itemize}
	\end{itemize}
	
\end{document}